\chapter{Equivariant KK-theory}
\label{ch:equivariant-kk}

\section{Hilbert modules and operator ideals}
In this section $B$ is a possibly nonunital complex $C^*$-algebra and
$E,F,G$ are right pre-Hilbert $B$-modules, with inner product linear in the
second variable. A Hilbert module is complete in the induced norm.
The first results below do not require completeness.

The pinned Mathlib scalar convention is implemented using
\texttt{CStarModule} over the opposite algebra $B^{\mathrm{op}}$:
$x b$ is represented by $\operatorname{op}(b)\mathbin{\bullet}x$, and
our $B$-valued inner product is the unop of Mathlib's inner product.
This fixes the multiplication order for the later Kasparov-cycle construction.

\begin{proposition}\label{prop:module-inner-separation}
\lean{BC4lean.KKTheory.module_inner_ext_left,BC4lean.KKTheory.module_inner_ext_right}
\leanok
The inner product separates points in either variable. In particular,
if $\langle z,x\rangle=\langle z,y\rangle$ for all $z\in E$, then $x=y$;
the same holds when the test vector is in the second variable.
\end{proposition}
\begin{proof}
\leanok
Apply definiteness to $\langle x-y,x-y\rangle$ and use conjugate symmetry.
\end{proof}

\begin{proposition}\label{prop:module-action-bound}
\uses{prop:module-inner-separation}
\lean{BC4lean.KKTheory.module_add_smul,BC4lean.KKTheory.module_smul_add,BC4lean.KKTheory.module_mul_smul,BC4lean.KKTheory.module_complex_smul,BC4lean.KKTheory.module_norm_smul_le}
\leanok
The coefficient action satisfies
$x(b+d)=xb+xd$, $(x+y)b=xb+yb$, $(xb)d=x(bd)$,
$x(cb)=c(xb)$ for $c\in\mathbb C$, and
$\|xb\|\leq\|x\|\|b\|$. These properties follow from
the module-valued inner-product identities and its norm, rather than being
additional assumptions on the action.
\end{proposition}
\begin{proof}
\leanok
Use separation by inner products for the algebraic identities. For the bound,
apply submultiplicativity to
$\langle xb,xb\rangle=b^*\langle x,x\rangle b$ and take square roots.
\end{proof}

\begin{definition}\label{def:adjointable-module-map}
\lean{BC4lean.KKTheory.AdjointableMap}
\leanok
An adjointable map $T:E\to F$ is a bounded complex-linear map together with
a bounded complex-linear map $T^*:F\to E$ satisfying
$\langle Tx,y\rangle=\langle x,T^*y\rangle$ for all $x,y$.
\end{definition}

\begin{proposition}\label{prop:adjointable-module-map-laws}
\uses{def:adjointable-module-map,prop:module-inner-separation}
\lean{BC4lean.KKTheory.AdjointableMap.adjoint_unique,BC4lean.KKTheory.AdjointableMap.map_op_smul,BC4lean.KKTheory.AdjointableMap.adjoint_adjoint,BC4lean.KKTheory.AdjointableMap.adjoint_comp,BC4lean.KKTheory.AdjointableMap.id_comp,BC4lean.KKTheory.AdjointableMap.comp_id,BC4lean.KKTheory.AdjointableMap.comp_assoc,BC4lean.KKTheory.AdjointableMap.adjoint_add,BC4lean.KKTheory.AdjointableMap.adjoint_smul}
\leanok
The adjoint is unique, and $T$ is automatically right $B$-linear.
Adjointable maps are closed under addition, complex scalar multiplication
and composition. They have identity and zero maps, composition is associative,
and
\[
 (T^*)^*=T,\qquad (ST)^*=T^*S^*,\qquad
 (S+T)^*=S^*+T^*,\qquad (cT)^*=\overline c\,T^*.
\]
\end{proposition}
\begin{proof}
\leanok
Apply the adjoint identity and inner-product separation. Construct the
operations using bounded complex-linear maps and their indicated adjoints.
\end{proof}

\begin{definition}\label{def:module-rank-one}
\uses{prop:module-action-bound,def:adjointable-module-map}
\lean{BC4lean.KKTheory.moduleRankOneCLM,BC4lean.KKTheory.moduleRankOne}
\leanok
For $x\in F$ and $y\in E$, define the rank-one module operator
\[
 \theta_{x,y}:E\longrightarrow F,\qquad
 \theta_{x,y}(z)=x\langle y,z\rangle.
\]
This is a bounded adjointable map. ``Rank-one'' is in the Hilbert-module sense;
it need not have one-dimensional complex range.
\end{definition}

\begin{proposition}\label{prop:module-rank-one-laws}
\uses{def:module-rank-one,prop:adjointable-module-map-laws}
\lean{BC4lean.KKTheory.moduleRankOne_norm_le,BC4lean.KKTheory.moduleRankOne_adjoint,BC4lean.KKTheory.comp_moduleRankOne,BC4lean.KKTheory.moduleRankOne_comp,BC4lean.KKTheory.moduleRankOne_comp_moduleRankOne}
\leanok
For all appropriately typed vectors and adjointable maps,
\[
 \|\theta_{x,y}\|\leq\|x\|\|y\|,\qquad
 \theta_{x,y}^*=\theta_{y,x},\qquad
 T\theta_{x,y}=\theta_{Tx,y},\qquad
 \theta_{x,y}T=\theta_{x,T^*y}.
\]
In particular,
$\theta_{x,y}\theta_{u,v}=\theta_{x\langle y,u\rangle,v}$.
\end{proposition}
\begin{proof}
\leanok
Use the module-action bound and Cauchy--Schwarz for the norm estimate.
The remaining identities follow by evaluation on vectors and the adjoint identity.
\end{proof}

\begin{proposition}\label{prop:standard-module-rank-one}
\uses{def:module-rank-one}
\lean{BC4lean.KKTheory.standard_module_inner,BC4lean.KKTheory.standard_module_complete,BC4lean.KKTheory.standard_moduleRankOne}
\leanok
The standard right module $B_B$ is complete, has inner product
$\langle x,y\rangle=x^*y$, and its rank-one operator $\theta_{x,y}$ is left
multiplication by $xy^*$.
\end{proposition}
\begin{proof}
\leanok
Use the existing complete $C^*$-algebra norm and evaluate the inner product
and rank-one construction on the opposite-algebra model of $B_B$.
\end{proof}

\begin{proposition}\label{prop:adjointable-module-norms}
\uses{def:adjointable-module-map}
\lean{BC4lean.KKTheory.AdjointableMap.norm_adjoint_toCLM,BC4lean.KKTheory.AdjointableMap.norm_toCLM_adjoint_comp_self}
\leanok
Adjoints preserve the operator norm and satisfy $\|T^*T\|=\|T\|^2$,
also for maps between different right pre-Hilbert modules.
\end{proposition}
\begin{proof}\leanok
Use $\|Tx\|^2=\|\langle x,T^*Tx\rangle\|$, Cauchy--Schwarz and the
operator-norm bound, first for adjoint norm preservation and then for the
square-root bound in the $C^*$-identity.
\end{proof}

\begin{proposition}\label{prop:complete-adjointable-algebra}
\uses{prop:adjointable-module-norms,prop:adjointable-module-map-laws}
\lean{BC4lean.KKTheory.AdjointableMap.instNormedAddCommGroup,BC4lean.KKTheory.AdjointableMap.instNormedAlgebra,BC4lean.KKTheory.AdjointableMap.instStarRing,BC4lean.KKTheory.AdjointableMap.isClosed_range_toCLMPair,BC4lean.KKTheory.AdjointableMap.toCLMPair_isometry,BC4lean.KKTheory.AdjointableMap.instCompleteSpace,BC4lean.KKTheory.AdjointableMap.instCStarAlgebra}
\leanok
Adjointable maps carry the norm $\|T\|=\|T:E\to F\|$.
For complete $E,F$ this space is complete. Adjointable endomorphisms of
complete $E$ form a unital complex $C^*$-algebra, with multiplication by
composition and involution by adjoint.
\end{proposition}
\begin{proof}\leanok
The adjoint identity defines a closed subset of the product of the two
bounded-operator spaces. Pairing $T$ with $T^*$ is an isometry because
$\|T^*\|=\|T\|$. Transport completeness from that closed subset, and use
the already proved composition, scalar, involution and $C^*$-norm laws.
\end{proof}

\begin{definition}\label{def:compact-module-operators}
\uses{prop:complete-adjointable-algebra,def:module-rank-one}
\lean{BC4lean.KKTheory.moduleFiniteRank,BC4lean.KKTheory.moduleCompact,BC4lean.KKTheory.IsModuleCompact,BC4lean.KKTheory.mem_moduleCompact_iff_approx}
\leanok
The finite-rank module maps are the complex linear span of the $\theta_{x,y}$.
The compact module maps $\mathcal K_B(E,F)$ are its operator-norm closure.
Equivalently, for every $\epsilon>0$ a compact map has a finite-rank module
approximation within $\epsilon$. This definition does not use complex
finite-dimensionality of the range.
\end{definition}

\begin{proposition}\label{prop:compact-module-ideal}
\uses{def:compact-module-operators,prop:module-rank-one-laws}
\lean{BC4lean.KKTheory.moduleCompact_isClosed,BC4lean.KKTheory.moduleCompact_adjoint,BC4lean.KKTheory.comp_mem_moduleCompact,BC4lean.KKTheory.mem_moduleCompact_comp,BC4lean.KKTheory.compactModuleOperatorIdeal_isTwoSided,BC4lean.KKTheory.compactModuleOperator_nonUnitalCStarAlgebra}
\leanok
Compact module maps are closed under adjoints. Composition with an
adjointable map on either side preserves compactness. Thus
$\mathcal K_B(E)$ is a closed adjoint-stable two-sided ideal of
$\mathcal L_B(E)$, and is a possibly nonunital $C^*$-algebra for complete $E$.
\end{proposition}
\begin{proof}\leanok
The rank-one formulas give the assertions on the linear span. Boundedness
of composition and adjoint norm preservation extend them to its closure.
\end{proof}

\begin{definition}\label{def:module-gradings}
\uses{def:adjointable-module-map}
\lean{BC4lean.KKTheory.CStarGrading,BC4lean.KKTheory.GradedHilbertModule,BC4lean.KKTheory.GradedHilbertModule.IsEven,BC4lean.KKTheory.GradedHilbertModule.IsOdd}
\leanok
A coefficient grading is an involutive complex star-algebra automorphism
$\beta$ of $B$. A compatible module grading is an involutive complex-linear
isometry $J$ with $\langle Jx,Jy\rangle=\beta(\langle x,y\rangle)$.
An operator is even when it commutes with the grading and odd when it
anticommutes. For trivially graded coefficients, $J$ is a selfadjoint unitary
adjointable operator. Completeness of the module is imposed when using it as
a Hilbert module.
\end{definition}

\begin{definition}\label{def:compatible-module-actions}
\uses{def:module-gradings}
\lean{BC4lean.KKTheory.CStarAlgebraAction,BC4lean.KKTheory.EquivariantHilbertModule,BC4lean.KKTheory.EquivariantHilbertModule.map_op_smul,BC4lean.KKTheory.EquivariantHilbertModule.continuous_action,BC4lean.KKTheory.EquivariantHilbertModule.PreservesGrading}
\leanok
For a discrete group $\Gamma$ acting on $B$ by complex star-algebra
automorphisms $\alpha_g$, a compatible action on $E$ consists of
complex-linear isometries $U_g$ satisfying the group laws and
$\langle U_gx,U_gy\rangle=\alpha_g(\langle x,y\rangle)$.
Then $U_g(xb)=(U_gx)\alpha_g(b)$, and the action is jointly continuous.
A graded action also commutes with $J$. The individual $U_g$ are generally
semilinear over $B$.
\end{definition}

\begin{proposition}\label{prop:equivariant-compact-conjugation}
\uses{def:compatible-module-actions,prop:compact-module-ideal}
\lean{BC4lean.KKTheory.EquivariantHilbertModule.conjugateOperator,BC4lean.KKTheory.EquivariantHilbertModule.conjugateMap_mem_moduleCompact,BC4lean.KKTheory.EquivariantHilbertModule.conjugateOperator_mem_moduleCompact}
\leanok
Conjugation $T\mapsto U_gTU_{g^{-1}}$ preserves adjointability, its operator
norm, and compact module operators. This applies to compatible maps between
different modules as well as endomorphisms.
\end{proposition}
\begin{proof}\leanok
Transport the adjoint identity through the coefficient automorphism.
Isometries preserve the norm, and conjugation sends $\theta_{x,y}$ to
$\theta_{U_gx,U_gy}$, so it preserves the closed rank-one span.
\end{proof}

\begin{definition}\label{def:countably-generated-module}
\uses{def:hilbert-modules}
\lean{BC4lean.KKTheory.moduleGeneratingSpan,BC4lean.KKTheory.IsCountablyGeneratedModule,BC4lean.KKTheory.moduleGeneratingSpan_closure_eq_top_iff,BC4lean.KKTheory.IsCountablyGeneratedModule.map_linearIsometryEquiv}
\leanok
A right Hilbert module is countably generated if there is a sequence $\xi_n$
such that the complex linear span of the vectors $\xi_n b$, $b\in B$, is norm
dense. This property is preserved by coefficient-compatible complex-linear
isometric module isomorphisms, including semilinear coefficient twists.
\end{definition}

\begin{proposition}\label{prop:standard-module-compact-algebra}
\uses{prop:standard-module-rank-one,prop:compact-module-ideal}
\lean{BC4lean.KKTheory.standardLeftMul,BC4lean.KKTheory.norm_standardLeftMul,BC4lean.KKTheory.moduleCompact_standard_eq_range,BC4lean.KKTheory.standardCompactEquiv,BC4lean.KKTheory.norm_standardCompactEquiv}
\leanok
Left multiplication identifies $B$ isometrically as a complex star-algebra
with $\mathcal K_B(B_B)$, also for nonunital $B$.
\end{proposition}
\begin{proof}\leanok
Test left multiplication by $b$ on $b^*$ to obtain its norm. Its image is
closed and contains every rank-one module operator. An approximate unit
shows conversely that every left multiplication is in the compact closure.
\end{proof}

\begin{definition}\label{def:hilbert-modules}
\uses{prop:complete-adjointable-algebra,prop:compact-module-ideal,def:module-gradings,def:compatible-module-actions,prop:standard-module-compact-algebra}
\lean{BC4lean.KKTheory.GradedHilbertModule,BC4lean.KKTheory.EquivariantHilbertModule,BC4lean.KKTheory.AdjointableMap.instCStarAlgebra,BC4lean.KKTheory.CompactModuleOperator}
\leanok
The preceding constructions give graded right Hilbert $C^*$-modules,
with a complete module norm, their adjointable and compact operators, and
continuous compatible actions of discrete groups. Coefficient algebras may be
nonunital and may have a nontrivial grading. Module completeness is an explicit
hypothesis, and compatibility of actions with the grading is explicit.
\end{definition}

\section{Equivariant Kasparov cycles}
\begin{definition}\label{def:even-kasparov-cycle}
\uses{def:countably-generated-module,def:module-gradings,def:compatible-module-actions,def:compact-module-operators}
\lean{BC4lean.KKTheory.KasparovCycle,BC4lean.KKTheory.KasparovCycle.IsDegenerate,BC4lean.KKTheory.KasparovCycle.zeroRepresentation,BC4lean.KKTheory.KasparovCycle.zeroRepresentation_isDegenerate}
\leanok
For trivially graded complex possibly nonunital $A,B$ and a discrete group
$\Gamma$, an even equivariant Kasparov cycle consists of a complete,
countably generated graded right Hilbert $B$-module $E$, a compatible action
commuting with the grading, an even equivariant possibly degenerate
star representation $\phi:A\to\mathcal L_B(E)$ and an odd adjointable
operator $F$. For all $a\in A$ and $g\in\Gamma$ require
\[
 (F-F^*)\phi(a),\quad (F^2-1)\phi(a),\quad
 F\phi(a)-\phi(a)F,\quad (U_gFU_{g^{-1}}-F)\phi(a)
 \ \in\ \mathcal K_B(E).
\]
A degenerate cycle has all four expressions zero. A zero representation with
zero operator is such a cycle on every countably generated compatible graded
module. The group action on $F$ is required only modulo compactness in the
stated local sense.
\end{definition}

\begin{definition}\label{def:odd-kasparov-cycle-data}
\uses{def:even-kasparov-cycle}
\lean{BC4lean.KKTheory.OddKasparovCycle,BC4lean.KKTheory.OddKasparovCycle.IsDegenerate,BC4lean.KKTheory.OddKasparovCycle.zeroRepresentation,BC4lean.KKTheory.OddKasparovCycle.zeroRepresentation_isDegenerate}
\leanok
The ungraded Fredholm picture of odd cycles uses the same complete,
countably generated equivariant module and the same four compactness
conditions, with no module grading. Its equivalence to the Clifford-graded
picture, and the homotopy quotient in either parity, require further proofs.
\end{definition}

\begin{proposition}\label{prop:standard-identity-cycle}
\uses{def:even-kasparov-cycle,prop:standard-module-compact-algebra}
\lean{BC4lean.KKTheory.standard_module_isCountablyGenerated,BC4lean.KKTheory.standardIdentityCycle,BC4lean.KKTheory.standardIdentityCycle_operator,BC4lean.KKTheory.standardIdentityCycle_representation}
\leanok
For separable possibly nonunital $B$, the standard right module $B_B$ is
countably generated. With all-even grading, the coefficient action, left
multiplication representation and $F=0$, it is an actual even equivariant
Kasparov cycle. These are the standard identity-cycle data; product unit
laws require the product construction below.
\end{proposition}
\begin{proof}\leanok
A countable dense sequence and an approximate unit give countable generation.
The square defect is $-L_b$, compact by the standard-module identification,
and the other defects vanish. Coefficient automorphisms give equivariance.
\end{proof}

\begin{definition}\label{def:fixed-module-operator-homotopy}
\uses{def:even-kasparov-cycle}
\lean{BC4lean.KKTheory.KasparovCycle.OperatorConditions,BC4lean.KKTheory.KasparovCycle.withOperator,BC4lean.KKTheory.KasparovOperatorHomotopy,BC4lean.KKTheory.KasparovOperatorHomotopy.cycleAt,BC4lean.KKTheory.KasparovOperatorHomotopy.continuous_operator,BC4lean.KKTheory.KasparovOperatorHomotopy.source_cycleAt,BC4lean.KKTheory.KasparovOperatorHomotopy.target_cycleAt,BC4lean.KKTheory.KasparovOperatorHomotopy.refl,BC4lean.KKTheory.KasparovOperatorHomotopy.symm,BC4lean.KKTheory.KasparovOperatorHomotopy.trans}
\leanok
On a fixed module, grading, representation and action, an operator homotopy
is an operator-norm-continuous path satisfying oddness and all four compact
defect conditions at every parameter. It yields actual cycles throughout,
including both endpoints. Constant paths, reversal and concatenation preserve
these conditions. Equivalence to homotopies over $C([0,1],B)$ and the full
KK quotient are further targets.
\end{definition}

\begin{proposition}\label{prop:scalar-creation-map-caution}
\lean{BC4lean.KKTheory.scalarCreationMap_adjoint_identity,BC4lean.KKTheory.scalarCreationMap_one_not_compact,BC4lean.KKTheory.ScalarHilbertModule.isCompactOperator_of_isModuleCompact,BC4lean.KKTheory.scalarCreationModuleMap_one_not_moduleCompact,BC4lean.KKTheory.ScalarCreationCycleCounterexample.firstCycle,BC4lean.KKTheory.ScalarCreationCycleCounterexample.secondCycle,BC4lean.KKTheory.ScalarCreationCycleCounterexample.secondCycle_isDegenerate,BC4lean.KKTheory.ScalarCreationCycleCounterexample.fullCycleCounterexample}
\leanok
Scalar creation maps in $\mathbb C\otimes_{\mathbb C}H$ are bounded
adjointable maps, but are not automatically compact. In the scalar correspondence
$\mathbb C\otimes_{\mathbb C}\ell^2(\mathbb N)\cong\ell^2(\mathbb N)$,
creation at $1$ is the identity under this identification, hence is not a
compact Hilbert-space operator. The same actual tensor creation map is
also noncompact for the scalar Hilbert-module compactness predicate.
The failure persists for two actual complete, countably generated graded
Kasparov cycles: the first is $\mathbb C$ with operator zero, and the second
is $\ell^2(\mathbb N)\oplus\ell^2(\mathbb N)$ with diagonal grading and odd
selfadjoint unitary flip. The second cycle has exactly zero defects.
\end{proposition}
\begin{proof}\leanok
The scalar tensor isometric identification gives the bounded adjoint
explicitly. Composing creation at $1$ with that identification yields the
identity. Its compactness would force $\ell^2(\mathbb N)$ to be
finite-dimensional, contradicting its independent delta vectors.
\end{proof}

\begin{proposition}\label{prop:rectangular-module-compactness}
\uses{prop:compact-module-ideal}
\lean{BC4lean.KKTheory.AdjointableMap.norm_sub_comp_sq_le,BC4lean.KKTheory.isModuleCompact_iff_adjoint_comp_self,BC4lean.KKTheory.IsModuleCompact.of_adjoint_comp_self}
\leanok
For an adjointable map $T:E\to F$ with complete source module,
$T$ is compact if and only if $T^*T$ is compact. The target need not be
complete. An approximate unit of $\mathcal K_B(E)$ and the rectangular
$C^*$-identity show that $T$ is a norm limit of compact maps $Te$.
\end{proposition}

\begin{definition}\label{def:hilbert-module-unitary-equivalence}
\uses{def:hilbert-modules}
\lean{BC4lean.KKTheory.HilbertModuleEquiv,BC4lean.KKTheory.HilbertModuleEquiv.toAdjointable,BC4lean.KKTheory.HilbertModuleEquiv.transportStarAlgEquiv,BC4lean.KKTheory.HilbertModuleEquiv.transportMap_norm,BC4lean.KKTheory.HilbertModuleEquiv.transportMap_isModuleCompact_iff}
\leanok
A Hilbert-module unitary equivalence is an actual complex-linear isometry
preserving the coefficient-valued inner product. It is coefficient-linear,
its inverse is its adjoint, and conjugation preserves operator norm, adjoints
and compactness, including rectangular operators.
\end{definition}

\begin{proposition}\label{prop:hilbert-module-direct-sum}
\uses{def:hilbert-modules,def:countably-generated-module}
\lean{BC4lean.KKTheory.HilbertModuleSum,BC4lean.KKTheory.inner_sum,BC4lean.KKTheory.norm_sum,BC4lean.KKTheory.hilbertModuleSum_complete,BC4lean.KKTheory.sumMap,BC4lean.KKTheory.IsModuleCompact.sumMap,BC4lean.KKTheory.IsCountablyGeneratedModule.sum,BC4lean.KKTheory.sumSwap,BC4lean.KKTheory.sumAssoc}
\leanok
The Hilbert-module direct sum has inner product
$\langle(x,y),(x',y')\rangle=\langle x,x'\rangle+\langle y,y'\rangle$
and norm $\sqrt{\|\langle x,x\rangle+\langle y,y\rangle\|}$.
Complete summands give a complete sum. Coordinate maps are adjointable,
diagonal compact maps are compact, and countable generation is preserved.
Swapping and reassociating summands give coefficient-inner-preserving
isometries.
\end{proposition}

\begin{definition}\label{def:kasparov-cycle-unitary-isomorphism}
\uses{def:even-kasparov-cycle,def:hilbert-module-unitary-equivalence}
\lean{BC4lean.KKTheory.KasparovCycleIso,BC4lean.KKTheory.KasparovCycleIso.refl,BC4lean.KKTheory.KasparovCycleIso.symm,BC4lean.KKTheory.KasparovCycleIso.trans,BC4lean.KKTheory.KasparovCycle.transport,BC4lean.KKTheory.KasparovCycle.transportIso,BC4lean.KKTheory.KasparovCycle.transport_isDegenerate,BC4lean.KKTheory.BundledKasparovCycle,BC4lean.KKTheory.BundledKasparovCycle.isomorphismSetoid}
\leanok
Cycle isomorphisms are genuine Hilbert-module unitaries intertwining the
grading, action, representation and Fredholm operator. They give an
equivalence relation across actual module types in a fixed explicit universe.
Unitary transport constructs all target cycle data and preserves degeneracy.
This is the isomorphism relation; the full homotopy quotient is a separate
construction.
\end{definition}

\begin{proposition}\label{prop:kasparov-cycle-direct-sum}
\uses{prop:hilbert-module-direct-sum,def:even-kasparov-cycle,def:odd-kasparov-cycle-data}
\lean{BC4lean.KKTheory.KasparovCycle.directSum,BC4lean.KKTheory.KasparovCycle.directSum_isDegenerate,BC4lean.KKTheory.OddKasparovCycle.directSum,BC4lean.KKTheory.OddKasparovCycle.directSum_isDegenerate}
\leanok
Even and odd raw cycles have genuine direct sums, with diagonal
representations and Fredholm operators and the compatible summed actions.
Every localized compact defect is the diagonal sum of the original defects.
Direct sums of degenerate cycles are degenerate. Group laws on a homotopy
quotient remain a further target.
\end{proposition}

\begin{proposition}\label{prop:compact-perturbation-operator-homotopy}
\uses{def:fixed-module-operator-homotopy}
\lean{BC4lean.KKTheory.affineKasparovOperator_square_defect,BC4lean.KKTheory.KasparovCycle.operatorConditions_affine_compactPerturbation,BC4lean.KKTheory.KasparovOperatorHomotopy.compactPerturbation}
\leanok
Two admissible operators $F_0,F_1$ on the same cycle frame give an actual
norm-continuous affine operator homotopy if
$(F_1-F_0)\phi(a)$ is compact for every $a$. The exact noncommutative square
identity controls the additional term $(F_1-F_0)^2\phi(a)$; commutation of the
operators and global compactness of their difference are not required.
\end{proposition}

\begin{proposition}\label{prop:continuous-section-hilbert-module}
\uses{def:hilbert-modules,def:countably-generated-module,prop:compact-module-ideal}
\lean{BC4lean.KKTheory.exists_finite_partition_approximation,BC4lean.KKTheory.continuousSectionCStarModule,BC4lean.KKTheory.continuousSectionOperator,BC4lean.KKTheory.continuousSectionOperator_isModuleCompact,BC4lean.KKTheory.continuousSection_isCountablyGenerated,BC4lean.KKTheory.continuousSectionFibreIsometryEquiv,BC4lean.KKTheory.continuousSectionFibreModuleEquiv}
\leanok
For compact Hausdorff $X$, continuous sections $C(X,E)$ with their supremum
norm form a right Hilbert $C(X,B)$-module under pointwise action and inner
product. Norm-continuous compact operator fields lift to compact module
operators by uniform finite partition approximation. Constant sections of a
countable generating family generate densely. The evaluation-kernel quotient
has its proved quotient norm and is unitarily identified with $E$.
This fibre assertion concerns continuous-section modules; fibres of arbitrary
$C(X,B)$-modules require a further construction.
\end{proposition}

\begin{definition}\label{def:balanced-module-tensor-algebraic}
\uses{def:hilbert-modules}
\lean{BC4lean.KKTheory.BalancedModuleTensor,BC4lean.KKTheory.moduleTensorMk_balance,BC4lean.KKTheory.moduleTensorLift,BC4lean.KKTheory.moduleTensorLift_unique,BC4lean.KKTheory.moduleTensorPureInner,BC4lean.KKTheory.moduleTensorPureInner_self_nonneg,BC4lean.KKTheory.moduleTensorPureInner_norm_le,BC4lean.KKTheory.moduleTensorInner,BC4lean.KKTheory.moduleTensorInner_star}
\leanok
Given an actual left action $\Psi:B\to\mathcal L_C(F)$, quotient the complex
algebraic tensor by the span of $xb\otimes y-x\otimes\Psi(b)y$.
This constructed quotient has the genuine balanced universal property.
Its coefficient pairing descends to a Hermitian form on the constructed quotient.
The elementary formula is
$\langle y,\Psi(\langle x,x'\rangle)y'\rangle$, respects both balance
relations, is positive on elementary vectors and satisfies the four-factor
norm bound. Separation and norm completion are additional obligations.
\end{definition}

\begin{proposition}\label{prop:injective-star-hom-positivity}
\uses{def:hilbert-modules}
\lean{BC4lean.KKTheory.injectiveStarHom_isSelfAdjoint_iff,BC4lean.KKTheory.injectiveStarHom_nonneg_iff,BC4lean.KKTheory.injectiveStarHom_le_iff}
\leanok
An injective nonunital complex star homomorphism between $C^*$-algebras
reflects selfadjointness and the given star-compatible positive order.
Functoriality of the continuous functional calculus for the negative part
proves positivity reflection.
\end{proposition}

\begin{proposition}\label{prop:hilbert-module-gram-positivity}
\uses{prop:hilbert-module-direct-sum,prop:injective-star-hom-positivity}
\lean{BC4lean.KKTheory.FiniteStandardModule,BC4lean.KKTheory.moduleGramSynthesis,BC4lean.KKTheory.moduleMatrixRepresentation,BC4lean.KKTheory.moduleMatrixRepresentation_injective,BC4lean.KKTheory.moduleGramMatrix_nonneg}
\leanok
For any finite family in a complete Hilbert $B$-module, its matrix of
coefficient inner products is positive in the actual $C^*$-matrix algebra
over $B$. Finite synthesis has the adjoint given by coefficient analysis.
A faithful matrix action and its linking-module embedding realize the Gram
matrix as $T^*T$ and reflect positivity. Nonunital $B$ is allowed.
\end{proposition}

\begin{proposition}\label{prop:hilbert-module-completion}
\uses{def:hilbert-modules,def:countably-generated-module}
\lean{BC4lean.KKTheory.completionModuleSMul,BC4lean.KKTheory.completionCStarModule,BC4lean.KKTheory.completionInner_coe,BC4lean.KKTheory.completion_norm_eq_sqrt_inner,BC4lean.KKTheory.hilbertModuleToCompletion,BC4lean.KKTheory.hilbertModuleToCompletion_denseRange,BC4lean.KKTheory.IsCountablyGeneratedModule.completion,BC4lean.KKTheory.AdjointableMap.completion,BC4lean.KKTheory.AdjointableMap.norm_completion,BC4lean.KKTheory.AdjointableMap.completionHom}
\leanok
The actual norm completion of a right pre-Hilbert $B$-module carries its
extended coefficient action and inner product, satisfies every Hilbert-module
law and retains the inherited completion norm. The canonical inclusion is a
dense coefficient-inner-preserving linear isometry. Countable generation is
preserved. Adjointable maps extend uniquely to completions, preserving
adjoints, norm, composition and the endomorphism star-algebra structure.
Completeness of the original module and a coefficient unit are not required.
\end{proposition}

\begin{proposition}\label{prop:balanced-module-tensor-positivity}
\uses{def:balanced-module-tensor-algebraic,prop:hilbert-module-gram-positivity}
\lean{BC4lean.KKTheory.moduleMatrixQuadratic_nonneg,BC4lean.KKTheory.moduleMatrixPairing_eq_zero_of_quadratic,BC4lean.KKTheory.moduleTensorInner_sum_mk_self_nonneg,BC4lean.KKTheory.moduleTensorInner_self_nonneg}
\leanok
The actual balanced tensor form is positive on every algebraic tensor, for
complete Hilbert modules. Amplifying the positive finite Gram matrix by the
left representation and using its genuine matrix positive cone proves
positivity on finite tensor sums. A zero quadratic vector pairs to zero with
every vector in the same positive matrix form. This uses actual second-module
inner products, with no assumed tensor norm or separation axiom.
\end{proposition}

\begin{proposition}\label{prop:nonunital-cstar-hom-range}
\lean{BC4lean.KKTheory.nonUnitalCStarHom_exists_norm_eq_lift,BC4lean.KKTheory.nonUnitalCStarHom_quotient_norm_mk,BC4lean.KKTheory.nonUnitalCStarHomQuotientIsometry,BC4lean.KKTheory.nonUnitalCStarHom_isClosed_range}
\leanok
Every nonunital complex star homomorphism of complete $C^*$-algebras has
closed range. The norm on its actual kernel quotient equals the image norm,
and every element in the image has a lift with exactly its image norm.
Continuous functional calculus trims a representative without changing its
image and gives the quotient isometry.
\end{proposition}

\begin{proposition}\label{prop:coefficient-au-on-modules}
\uses{prop:module-action-bound}
\lean{BC4lean.KKTheory.tendsto_coefficient_approximateUnit_smul,BC4lean.KKTheory.exists_positive_selfadjoint_contraction_smul_approx_on}
\leanok
The actual increasing positive contractive coefficient approximate unit acts
strongly as the identity on every right Hilbert module. A finite family of
vectors can be approximated simultaneously by one positive selfadjoint
contractive coefficient. A unit, completeness of the module and an extra
essentiality hypothesis are not required.
\end{proposition}

\begin{proposition}\label{prop:positive-module-form-quotient}
\uses{prop:module-action-bound}
\lean{BC4lean.KKTheory.PositiveModuleForm,BC4lean.KKTheory.PositiveModuleForm.inner_mul_inner_swap_le,BC4lean.KKTheory.PositiveModuleForm.formNorm_triangle,BC4lean.KKTheory.PositiveModuleForm.nullRadical,BC4lean.KKTheory.PositiveModuleQuotient,BC4lean.KKTheory.PositiveModuleQuotient.instCStarModule}
\leanok
A positive semidefinite coefficient-valued module form satisfies
Cauchy--Schwarz and gives the seminorm
$\|x\|_q=\sqrt{\|\langle x,x\rangle\|}$. Its null vectors form a
coefficient-invariant complex submodule. Quotienting that actual radical,
with this induced form norm, constructs a separated pre-Hilbert module.
No definiteness or pre-existing vector norm is assumed.
\end{proposition}

\begin{proposition}\label{prop:completed-interior-module-tensor}
\uses{def:balanced-module-tensor-algebraic,prop:balanced-module-tensor-positivity,prop:hilbert-module-completion}
\lean{BC4lean.KKTheory.moduleTensorInner_eq_zero_of_self,BC4lean.KKTheory.moduleTensorNull,BC4lean.KKTheory.preModuleTensorCStarModule,BC4lean.KKTheory.PreHilbertModuleTensor,BC4lean.KKTheory.InteriorModuleTensor,BC4lean.KKTheory.interiorModuleTensorInner_mk_mk,BC4lean.KKTheory.interiorModuleTensorMk_balance,BC4lean.KKTheory.interiorModuleTensor_span_dense}
\leanok
For complete Hilbert modules and an actual possibly degenerate nonunital
left representation, the tensor form's zero quadratic vectors are precisely
its pairing radical. The coefficient action descends through that radical.
The induced norm on the separated balancing quotient and its actual norm
completion construct $E\otimes_B F$. Elementary tensors have the prescribed
inner product, obey balance and span a dense subspace.
\end{proposition}

\begin{proposition}\label{prop:interior-tensor-operators}
\uses{prop:completed-interior-module-tensor,prop:complete-adjointable-algebra}
\lean{BC4lean.KKTheory.interiorModuleTensorOperator,BC4lean.KKTheory.interiorModuleTensorOperator_mk,BC4lean.KKTheory.interiorModuleTensorOperator_norm_le,BC4lean.KKTheory.interiorModuleTensorOperatorHom,BC4lean.KKTheory.interiorModuleTensorOperator_adjoint}
\leanok
An adjointable coefficient-linear $T:E\to E$ induces the actual
$T\otimes1$ on the completed interior tensor. It maps
$x\otimes y$ to $Tx\otimes y$, has norm at most $\|T\|$, and
preserves adjoints. The construction is a unital complex star-algebra
homomorphism on the adjointable endomorphism algebra.
\end{proposition}

\begin{proposition}\label{prop:standard-tensor-essential-range}
\uses{prop:completed-interior-module-tensor,prop:standard-module-compact-algebra}
\lean{BC4lean.KKTheory.moduleTensorInner_standard_eq,BC4lean.KKTheory.standardInteriorModuleTensorEvaluationCLM_norm,BC4lean.KKTheory.standardInteriorModuleTensorEvaluation_range,BC4lean.KKTheory.standardInteriorModuleTensorEquivEssential,BC4lean.KKTheory.standardInteriorModuleTensorHilbertEquiv}
\leanok
For a possibly degenerate left representation $\Psi:B\to\mathcal L_C(F)$,
the actual map $b\otimes y\mapsto\Psi(b)y$ identifies $B_B\otimes_B F$
unitarily with the closed complex span of $\Psi(B)F$. Surjectivity onto
all of $F$ is not asserted for a degenerate representation, and no
complement for this essential range is assumed.
\end{proposition}

\begin{definition}\label{def:general-graded-kasparov-cycle}
\uses{def:even-kasparov-cycle,def:module-gradings}
\lean{BC4lean.KKTheory.GradedKasparovCycle,BC4lean.KKTheory.GradedKasparovCycle.trivialEquiv,BC4lean.KKTheory.GradedKasparovCycleIso,BC4lean.KKTheory.GradedKasparovCycleIso.symm,BC4lean.KKTheory.GradedKasparovCycleIso.trans}
\leanok
For compatible algebra gradings $\delta_A,\delta_B$, a genuine graded
cycle uses a $\delta_B$-semilinear module grading, an equivariant graded
representation, an odd Fredholm operator, and the same localized compact
defects with commutator
$F\phi(a)-\phi(\delta_A(a))F$. Algebra actions preserve their gradings.
The trivial-grading specialization is exactly the earlier even cycle
definition. Actual whole-cycle unitaries give reflexive, symmetric and
transitive isomorphisms across module carriers.
\end{definition}

\begin{proposition}\label{prop:complex-clifford-one}
\lean{BC4lean.KKTheory.CliffordOne.grading,BC4lean.KKTheory.CliffordOne.epsilon_square,BC4lean.KKTheory.CliffordOne.epsilon_odd,BC4lean.KKTheory.CliffordOne.existsUnique_lift}
\leanok
The complex one-generator Clifford algebra is the actual $C^*$-algebra
$\mathbb C\times\mathbb C$ with grading exchanging the coordinates.
Its odd generator $\varepsilon=(1,-1)$ is a selfadjoint unitary.
Every selfadjoint unitary in a unital complex $C^*$-algebra gives a unique
unital star homomorphism sending $\varepsilon$ to that unitary.
\end{proposition}

\begin{proposition}\label{prop:clifford-coefficient-tensor}
\uses{prop:complex-clifford-one}
\lean{BC4lean.KKTheory.cliffordTensorCoordinates_tmul,BC4lean.KKTheory.cliffordTensorNorm_cstar,BC4lean.KKTheory.cliffordTensorCoordinateStarAlgEquiv,BC4lean.KKTheory.CliffordCommutingRepresentation.maximalNorm_eq}
\leanok
The actual algebraic tensor $B\otimes_{\mathbb C}\mathrm{Cl}_1$ has
coordinate map $b\otimes(z_+,z_-)\mapsto(z_+b,z_-b)$ and norm
$\max(\|b_+\|,\|b_-\|)$. Its constructed complete $C^*$-algebra is
star-isomorphic to $B\times B$. The norm is the maximal norm over commuting
representations, including nonunital and degenerate coefficient
representations. No spatial tensor-norm comparison is claimed here.
\end{proposition}

\begin{proposition}\label{prop:cycle-sum-unitaries-and-inverse-rotation}
\uses{prop:kasparov-cycle-direct-sum,def:fixed-module-operator-homotopy}
\lean{BC4lean.KKTheory.KasparovCycleIso.directSum,BC4lean.KKTheory.KasparovCycleIso.sumComm,BC4lean.KKTheory.KasparovCycleIso.sumAssoc,BC4lean.KKTheory.KasparovCycle.opposite,BC4lean.KKTheory.KasparovOperatorHomotopy.inverseRotation,BC4lean.KKTheory.KasparovOperatorHomotopy.inverseRotation_target_isDegenerate,BC4lean.KKTheory.KasparovOperatorHomotopy.directSum}
\leanok
Actual summed whole-cycle unitaries give exchange and associativity of
even raw direct sums. A cycle's opposite reverses its grading and Fredholm
operator. The explicit cosine--sine rotation on the sum of a cycle and its
opposite is an admissible norm-continuous operator homotopy ending at
the exactly degenerate flip cycle. Operator homotopies can be summed.
These statements do not themselves construct the full homotopy quotient.
\end{proposition}

\begin{proposition}\label{prop:operator-path-as-interval-cycle}
\uses{def:fixed-module-operator-homotopy,prop:continuous-section-hilbert-module}
\lean{BC4lean.KKTheory.CStarAlgebraAction.continuousSections,BC4lean.KKTheory.KasparovOperatorHomotopy.intervalCycle,BC4lean.KKTheory.KasparovOperatorHomotopy.fibreCycleIso,BC4lean.KKTheory.KasparovOperatorHomotopy.fibreCycleZeroIso,BC4lean.KKTheory.KasparovOperatorHomotopy.fibreCycleOneIso}
\leanok
An admissible fixed-module operator path gives an actual equivariant cycle
on the continuous-section Hilbert $C([0,1],B)$-module. Its evaluation-kernel
fibres carry the induced whole cycles, and genuine cycle unitaries identify
them with the original parameterwise cycles, including the endpoints.
This does not assert that every interval cycle has a fixed-module path model.
\end{proposition}

\begin{proposition}\label{prop:interior-tensor-creation}
\uses{prop:completed-interior-module-tensor,prop:rectangular-module-compactness}
\lean{BC4lean.KKTheory.moduleCreationMap,BC4lean.KKTheory.moduleCreationMap_adjoint_mk,BC4lean.KKTheory.moduleCreationMap_norm_le,BC4lean.KKTheory.moduleCreationMap_adjoint_comp,BC4lean.KKTheory.moduleCreationMap_isModuleCompact_iff,BC4lean.KKTheory.moduleCreationMap_isModuleCompact_of_compact_left_action}
\leanok
Creation $\theta_\xi(y)=\xi\otimes y$ on the actual completed tensor is
adjointable, has norm at most $\|\xi\|$, and satisfies
$\theta_\xi^*(\zeta\otimes y)=\Psi(\langle\xi,\zeta\rangle)y$ and
$\theta_\xi^*\theta_\zeta=\Psi(\langle\xi,\zeta\rangle)$.
Precisely, $\theta_\xi$ is compact if and only if
$\Psi(\langle\xi,\xi\rangle)$ is compact. Compactness of the entire left
action is sufficient; arbitrary creation maps are not asserted compact.
\end{proposition}

\begin{proposition}\label{prop:interior-tensor-countable-generation}
\uses{prop:completed-interior-module-tensor,prop:interior-tensor-creation,def:countably-generated-module}
\lean{BC4lean.KKTheory.IsCountablyGeneratedModule.interiorTensor}
\leanok
The actual completed interior tensor of two countably generated Hilbert
modules is countably generated. Pairing the two generating sequences and
using the bounded tensor maps proves density. Nonunital coefficients and
a degenerate left representation are allowed.
\end{proposition}

\begin{definition}\label{def:actual-kasparov-connection}
\uses{prop:interior-tensor-creation,def:even-kasparov-cycle}
\lean{BC4lean.KKTheory.KasparovConnection,BC4lean.KKTheory.KasparovConnection.add_compact}
\leanok
The analytic $S$-connection predicate on the actual tensor records both the operator and
adjoint compactness conditions
$\theta_\xi S-(-1)^{|\xi|}R\theta_\xi\in\mathcal K$
for every homogeneous $\xi$, with the analogous condition for $S^*,R^*$.
A globally compact perturbation of $R$ preserves all four compactness conditions.
The odd degree of a graded connection must be proved separately for the actual
tensor grading; preserving that degree also requires an odd perturbation.
Existence of a graded connection or a Kasparov product is not assumed.
\end{definition}

\begin{proposition}\label{prop:actual-interior-creation-counterexample}
\uses{prop:scalar-creation-map-caution,prop:interior-tensor-creation}
\lean{BC4lean.KKTheory.InteriorCreationCycleCounterexample.creationAtOne_adjoint_comp,BC4lean.KKTheory.InteriorCreationCycleCounterexample.creationAtOne_not_moduleCompact,BC4lean.KKTheory.InteriorCreationCycleCounterexample.fullCycleCounterexample}
\leanok
For the two full scalar input cycles described above, creation at $1$
into the constructed general interior tensor is noncompact. Its actual
adjoint square is the identity on the infinite-dimensional second module.
This verifies the counterexample within the constructed tensor, without
substituting a different tensor model.
\end{proposition}

\begin{proposition}\label{prop:gram-coefficient-quadratic-positivity}
\uses{prop:module-action-bound}
\lean{BC4lean.KKTheory.moduleGram_quadratic_eq_inner,BC4lean.KKTheory.moduleGram_quadratic_nonneg}
\leanok
Every finite Gram coefficient quadratic expression
$\sum_{i,j} b_i^*\langle x_i,x_j\rangle b_j$ equals
$\langle\sum_i x_i b_i,\sum_j x_j b_j\rangle$ and is positive.
This follows directly from the coefficient inner-product identities and
does not assume a tensor norm or tensor positivity.
\end{proposition}

\begin{proposition}\label{prop:arbitrary-hilbert-module-fibres}
\uses{prop:positive-module-form-quotient,prop:hilbert-module-completion}
\lean{BC4lean.KKTheory.HilbertModuleFibre,BC4lean.KKTheory.hilbertModuleFibre_inner_mk,BC4lean.KKTheory.hilbertModuleFibreMk_denseRange,BC4lean.KKTheory.IsCountablyGeneratedModule.fibre,BC4lean.KKTheory.hilbertModuleFibreOperatorHom,BC4lean.KKTheory.IsModuleCompact.fibre}
\leanok
For an arbitrary Hilbert $C(X,B)$-module over compact Hausdorff $X$,
evaluate its actual inner product at $t$, quotient its null radical,
and complete the resulting pre-Hilbert module. This constructs the fibre
$E_t$ and its canonical dense map, with the prescribed evaluated inner
product. Countable generation descends. Adjointable operators induce a
contractive unital star homomorphism on fibres, and rank-one evaluation
proves that compact module operators descend to compact operators.
\end{proposition}

\begin{proposition}\label{prop:arbitrary-module-fibre-quotient-norm}
\uses{prop:arbitrary-hilbert-module-fibres,prop:coefficient-au-on-modules}
\lean{BC4lean.KKTheory.hilbertModuleBanachFibre_norm_mk,BC4lean.KKTheory.hilbertModuleBanachFibreEquiv,BC4lean.KKTheory.evaluatedModuleQuotient_completeSpace,BC4lean.KKTheory.hilbertModuleFibreMk_surjective}
\leanok
For complete $E$ over compact Hausdorff $X$, its actual closed evaluated
radical quotient has Banach quotient norm equal to the evaluated inner-product
norm. Continuous cutoffs and the coefficient approximate unit give arbitrarily
close norm-controlled representatives. Consequently the evaluated quotient
is already complete and the canonical map $E\to E_t$ is surjective.
\end{proposition}


\begin{proposition}\label{prop:compact-au-acts-strongly}
\uses{prop:compact-module-ideal,prop:module-rank-one-laws}
\lean{BC4lean.KKTheory.moduleRankOne_self_norm,BC4lean.KKTheory.compact_approximateUnit_vector_error_sq_le,BC4lean.KKTheory.tendsto_canonical_compact_approximateUnit_apply,BC4lean.KKTheory.tendsto_compact_sequence_apply}
\leanok
The exact identity $\|\theta_{x,x}\|=\|x\|^2$ implies that the actual
compact-operator approximate unit acts strongly as the identity on $E$.
The residual rank-one compression bounds vector errors by compact
approximate-unit errors. This uses the original $B$-valued inner product;
it does not retype $E$ as a module with a compact-operator-valued inner product.
\end{proposition}

\begin{proposition}\label{prop:arbitrary-interval-cycle-fibres}
\uses{prop:arbitrary-hilbert-module-fibres,def:even-kasparov-cycle}
\lean{BC4lean.KKTheory.GradedHilbertModule.fibre,BC4lean.KKTheory.EquivariantHilbertModule.fibre,BC4lean.KKTheory.KasparovCycle.fibre,BC4lean.KKTheory.KasparovCycle.isDegenerate_fibre,BC4lean.KKTheory.KasparovOperatorHomotopy.intervalFibreIso}
\leanok
An actual equivariant Kasparov cycle on an arbitrary complete countably
generated Hilbert $C(X,B)$-module, with coefficient action acting pointwise
and fixing $X$, induces a whole cycle on every evaluated fibre.
Its grading, semilinear action, representation and Fredholm operator
are the induced actual fibre maps, and all four compact defects descend.
Degeneracy is preserved. For a fixed-module operator path, these general
fibres are unitarily identified with its original parameterwise cycles.
\end{proposition}

\begin{proposition}\label{prop:compact-fibre-epimorphism-and-kernel}
\uses{prop:arbitrary-module-fibre-quotient-norm,prop:nonunital-cstar-hom-range,prop:compact-au-acts-strongly}
\lean{BC4lean.KKTheory.hilbertModuleCompactFibreHom,BC4lean.KKTheory.hilbertModuleCompactFibreHom_surjective,BC4lean.KKTheory.exists_compactFibre_lift_norm_eq,BC4lean.KKTheory.compact_fibreKernel_eq_fibreNullCompact}
\leanok
Evaluation gives a surjective nonunital star homomorphism
$\mathcal K_{C(X,B)}(E)\to\mathcal K_B(E_t)$, with norm-preserving lifts.
Its kernel is exactly the closed complex span of the rank-one maps
$\theta_{x,y}$ with $x_t=0$. Fibre-vector surjectivity lifts rank-one maps;
closed $C^*$-homomorphism range gives the full epimorphism, and the actual
compact-algebra approximate unit proves the kernel description.
\end{proposition}

\begin{proposition}\label{prop:cone-hilbert-module}
\uses{prop:continuous-section-hilbert-module,prop:coefficient-au-on-modules}
\lean{BC4lean.KKTheory.ConeHilbertModule,BC4lean.KKTheory.ConeHilbertModule.instCStarModule,BC4lean.KKTheory.coneRampMap_denseRange,BC4lean.KKTheory.cone_isCountablyGenerated,BC4lean.KKTheory.coneOperatorHom}
\leanok
The actual closed section space $\{x\in C([0,1],E):x(0)=0\}$ is a complete
Hilbert $C([0,1],B)$-module with pointwise form and action. Uniform ramp
regularization proves the density needed for its explicit countable
generators when $E$ is countably generated. Constant adjointable operators
act on this cone and give a norm-contracting unital star homomorphism.
No general closed-submodule countable-generation assertion is used.
\end{proposition}


\begin{definition}\label{def:equivariant-kk}
\uses{def:hilbert-modules}
\notready
For separable $\Gamma$-$C^*$-algebras $A,B$, construct $KK_i^\Gamma(A,B)$
from graded equivariant Kasparov cycles and homotopy. The raw even and ungraded odd cycle conditions are specified above.
Construct a homotopy quotient with direct sum and inverse laws, and prove the
parity identification before claiming the resulting abelian groups. For the
intended applications, $A,B$ are separable and $\Gamma$ is countable discrete.
\end{definition}

\section{Kasparov products and functoriality}
\begin{proposition}\label{prop:kasparov-product}
\uses{def:equivariant-kk}
\notready
Construct the Kasparov product, with existence, well-definedness,
associativity, and units. Record contravariance in the first algebra and
covariance in the second. This is a substantial development, to be split
into supporting lemmas in the detailed blueprint.
\end{proposition}

\section{Comparison with operator K-theory}
\begin{proposition}\label{prop:kk-k-theory}
\uses{def:equivariant-kk,def:operator-k-theory}
\notready
For the trivial group, establish $KK_i(\mathbb C,B)\cong K_i(B)$.
\end{proposition}

General algebras in KK-theory are infrastructure for the construction;
the main conjecture of this project still has no coefficients.
